%==============================================================
%  EE201  Circuit Theory I  --  Fall 2026/27
%  Common syllabus (departmental) + Section 2 specific additions
%  Instructor of Section 2: Mert Ankaralı
%
%  Compile:  pdflatex 2026_Syllabus.tex  (twice, for hyperref)
%==============================================================
\documentclass[10pt,a4paper]{article}

%-------------------- packages --------------------------------
\usepackage[T1]{fontenc}
\usepackage[utf8]{inputenc}
\usepackage[english]{babel}
\usepackage{lmodern}
\usepackage[margin=1.8cm,top=1.9cm,bottom=1.7cm,headheight=13pt,headsep=8pt,footskip=20pt]{geometry}
\usepackage{array}
\usepackage{booktabs}
\usepackage{enumitem}
\usepackage[table]{xcolor}
\usepackage{fancyhdr}
\usepackage{titlesec}
\usepackage[most]{tcolorbox}
\usepackage[colorlinks=true,allcolors=metured,breaklinks=true]{hyperref}

%-------------------- colours ---------------------------------
\definecolor{metured}{RGB}{155,27,48}   % METU maroon
\definecolor{metugray}{RGB}{88,89,91}
\definecolor{lightgray}{RGB}{242,242,242}

%-------------------- course-wide macros ----------------------
% Change these in ONE place if anything moves.
\newcommand{\coursecode}{EE201}
\newcommand{\coursename}{Circuit Theory I}
\newcommand{\term}{Fall 2026--2027}
\newcommand{\mysection}{Section 2}
\newcommand{\myname}{Mert Ankaralı}
\newcommand{\myoffice}{A-410/D-104}
\newcommand{\mymail}{mertan@metu.edu.tr}
% >>> Public repository that hosts the lecture notes of Section 2 <<<
% Repo: mertankarali/Lecture-Notes (public), course folder: METU-EE201
\newcommand{\notesrepo}{https://github.com/mertankarali/Lecture-Notes/tree/master/METU-EE201}
\newcommand{\notesrepotext}{github.com/mertankarali/Lecture-Notes}
\newcommand{\odtuclass}{https://odtuclass.metu.edu.tr}
\newcommand{\coordinator}{Prof. T. Engin Tuncer}
\newcommand{\coordinatormail}{etuncer@metu.edu.tr}
\newcommand{\tunanotes}{https://users.metu.edu.tr/etuna/ee201/}
\newcommand{\tunanotestext}{users.metu.edu.tr/etuna/ee201}

%-------------------- compact layout --------------------------
\pagestyle{fancy}
\fancyhf{}
\lhead{\scriptsize\textcolor{metugray}{\coursecode\ -- Syllabus}}
\rhead{\scriptsize\textcolor{metugray}{METU EEE, \term}}
\cfoot{\scriptsize\thepage}
\renewcommand{\headrulewidth}{0.4pt}

\titleformat{\section}{\normalsize\bfseries\color{metured}}{\thesection.}{0.5em}{}
\titleformat{\subsection}[runin]
  {\small\bfseries\color{metugray}}{}{0pt}{}[\,---\,]
\titlespacing*{\section}{0pt}{0.9ex plus .2ex}{0.35ex}
\titlespacing*{\subsection}{0pt}{0.7ex plus .1ex}{0.5em}

\setlist[itemize]{itemsep=0.5pt,topsep=1.5pt,parsep=0pt,partopsep=0pt,leftmargin=1.4em}
\setlist[enumerate]{itemsep=0.5pt,topsep=1.5pt,parsep=0pt,partopsep=0pt,leftmargin=1.6em}
\setlength{\parindent}{0pt}
\setlength{\parskip}{0.45ex plus .1ex}
\renewcommand{\arraystretch}{1.05}
\setlength{\tabcolsep}{5pt}

\newtcolorbox{notebox}[1][]{
  enhanced, colback=lightgray, colframe=metured,
  boxrule=0.7pt, left=7pt, right=7pt, top=4pt, bottom=4pt,
  before skip=4pt, after skip=4pt, arc=2pt, #1}

\newcolumntype{L}[1]{>{\raggedright\arraybackslash}p{#1}}

%==============================================================
\begin{document}

%-------------------- title block -----------------------------
\begin{center}
  {\small\scshape Middle East Technical University \quad\textbar\quad
   Department of Electrical and Electronics Engineering}\\[3pt]
  {\Large\bfseries\color{metured} \coursecode: \coursename}\\[2pt]
  {\small \term\ \quad\textbar\quad Course Syllabus}
\end{center}
\vspace{-0.3em}
{\color{metured}\hrule height 0.8pt}
\vspace{0.4em}

%==============================================================
\section{Instructor Information}

\begin{center}
\small
{\setlength{\tabcolsep}{4pt}\begin{tabular}{c L{3.6cm} c L{4.3cm} @{\hskip 0.8em} c L{3.2cm}}
\toprule
\textbf{Sec.} & \textbf{Instructor} & \textbf{Office} & \textbf{e-mail} & \textbf{Sec.} & \textbf{Instructor}\\
\midrule
1 & T. Engin Tuncer & E109  & \href{mailto:etuncer@metu.edu.tr}{etuncer@metu.edu.tr} & 4 & Emre Tuna (E113)\\
\rowcolor{lightgray}
\textbf{2} & \textbf{\myname} & \textbf{\myoffice} & \href{mailto:\mymail}{\mymail} & 5 & Yunus Aslan (C-103)\\
3 & Emre Tuna & E113 & \href{mailto:etuna@metu.edu.tr}{etuna@metu.edu.tr} & &\\
\bottomrule
\end{tabular}}
\end{center}
\vspace{-0.4em}

\textbf{Course coordinator:} \coordinator\
(\href{mailto:\coordinatormail}{\coordinatormail}).

\subsection*{Teaching assistants}
\textbf{\coursecode\ (theory):} Anıl Gülbeden (\href{mailto:gulbeden@metu.edu.tr}{gulbeden@metu.edu.tr}).\\
\textbf{\coursecode/EE213 --- laboratory:} Alp Gürsoy, Eda Özkaynar, Orhun Koç, Civan Serhat Çevik,
Elif Çağla Tay, Enfal Çabuk, Adem Utku Atasayar, Özgür Gülsuna, Alper Alka.

%==============================================================
\section{Textbooks, References and Other Resources}

\begin{enumerate}
  \item Alexander, C.~K., \& Sadiku, M.~N.~O. (2016). \emph{Fundamentals of Electric Circuits.}
        McGraw-Hill. \hfill \textit{(main textbook)}
  \item Chua, L.~O., Desoer, C.~A., \& Kuh, E.~S. (1987). \emph{Linear and Nonlinear Circuits.}
        McGraw-Hill. \hfill \textit{(reference)}
  \item Lecture notes and videos by Emre Tuna, \href{\tunanotes}{\texttt{\tunanotestext}}.
\end{enumerate}

\textbf{ODTÜCLASS} (\href{\odtuclass}{odtuclass.metu.edu.tr}) is the official channel for all
sections: announcements, the detailed course outline, grades, exam and laboratory information,
and the common course material listed above.

%==============================================================
\section{Course Materials for \mysection}

In addition to the common resources above, the lecture notes and supplementary material
\emph{specific to \mysection} are maintained in a public repository:

\begin{notebox}
\begin{center}
  \textbf{\mysection\ lecture notes and course material}\\[2pt]
  {\ttfamily\href{\notesrepo}{\notesrepotext}}\;\;$\rightarrow$\;\;{\ttfamily METU-EE201}
\end{center}
\vspace{-0.2em}
\scriptsize
Public: no account is required to read or download. Updated \emph{during} the semester,
typically before each lecture --- refresh it regularly rather than downloading once.
Corrections and typo reports (issues or e-mail) are very welcome.
\end{notebox}
\vspace{-0.2em}

\textbf{What is where:} \emph{repository} $\rightarrow$ \mysection\ lecture notes and
supplementary material; \emph{ODTÜCLASS} $\rightarrow$ everything official and everything
common to all sections.
Announcements are always made on ODTÜCLASS.

%==============================================================
\section{Grading}

\begin{center}
\small
\begin{tabular}{l c c c c c @{\hskip 1.4em} c}
\toprule
\textbf{Component} & Attendance & Midterm 1 & Midterm 2 & Final Exam & Laboratory & \textbf{Total}\\
\midrule
\textbf{Weight} & 5\,\% & 20\,\% & 20\,\% & 35\,\% & 20\,\% & \textbf{100\,\%}\\
\bottomrule
\end{tabular}
\end{center}
\vspace{-0.3em}

%==============================================================
\section{Attendance \normalfont\small\color{metugray}(\mysection)}

Attendance is worth 5\,\% of the grade, and attendance below 40\,\% means \textbf{NA}
(see Section~6). In \mysection\ it is collected as follows.

\begin{itemize}
  \item Attendance is taken with a \textbf{QR code through ODTÜCLASS}. Come with a device that
        can read it.
  \item It is taken \textbf{once per class meeting} --- once in a one-hour meeting and once in a
        two-hour meeting, \emph{not} once per lecture hour.
  \item It is normally opened \textbf{at the beginning} of the meeting. In a two-hour meeting,
        if I give a break, I may open it \textbf{at the break} instead.
  \item Meetings are not always run the same way: I sometimes teach two hours as a single block,
        sometimes finish early, and I do not always open attendance at the break.
  \item Therefore there is \textbf{no pattern to rely on}. Missing the first hour and arriving at
        the break --- or any similar strategy --- is entirely at your own risk: if attendance was
        taken while you were not in the room, that meeting is lost and it is not repeated.
\end{itemize}

%==============================================================
\section{\coursecode\ Final Examination Policy}

A student (i)~whose attendance is below 40\,\%, (ii)~missing any midterm examination without a
valid excuse, (iii)~missing any laboratory experiment without a valid excuse, or (iv)~missing the
Laboratory Final Exam without a valid excuse, will \textbf{not} be admitted to the final
examination and will receive the \textbf{NA} grade.

%==============================================================
\section{Make-up Policy}

A make-up compensates for a \emph{single} missing exam or missing laboratory. Students missing
more than one exam or laboratory have to take multiple make-ups. Make-up exams can take place
before or after the final exam; please follow the announcements on \href{\odtuclass}{ODTÜCLASS}.
The laboratory make-up is announced separately.

%==============================================================
\section{Previous Students of \coursecode\ and EE213}

\begin{itemize}
  \item Students who have taken \textbf{EE213}, have a grade for it in their transcript, and are
        registered to the new \coursecode\ will get their course grade from the \coursecode\
        theoretical part (100\,\%); previous laboratory grades will \emph{not} be included.
  \item Students who \textbf{failed} EE213 and take only the EE213 laboratory will take their
        grades from the laboratory work only, as before.
  \item The new \coursecode\ (which includes laboratory work) is opened only for new students
        who take this course for the first time.
  \item If you are not sure which case applies to you, send an e-mail to the course coordinator,
        \coordinator\ (\href{mailto:\coordinatormail}{\coordinatormail}).
\end{itemize}

%==============================================================
\section{Course Outline}

\begin{enumerate}[leftmargin=1.5em,labelsep=0.4em]
  \item \textbf{Basic Concepts} \hfill \textit{(6 hrs)}\\
        Introduction; lumped elements and lumped circuits; interconnection equations;
        branch relations; basic lumped elements; circuit analysis.
  \item \textbf{Linear Time-Invariant Resistive Circuits} \hfill \textit{(10 hrs)}\\
        LTI resistive elements (resistors, dependent sources, ideal transformers);
        analysis methods; one-port circuits; two-ports.
  \item \textbf{Nonlinear Resistive Circuits} \hfill \textit{(3 hrs)}\\
        Nonlinear resistive elements and circuits; load-line analysis;
        piecewise-linear resistive circuits.
  \item \textbf{Operational Amplifier Circuits} \hfill \textit{(7 hrs)}\\
        Operational amplifiers; basic op-amp circuits; more realistic op-amp models;
        miscellaneous resistive op-amp circuits; op-amp with piecewise-linear elements.
  \item \textbf{Dynamic Elements} \hfill \textit{(5 hrs)}\\
        Ramp, step and impulse functions; capacitors; inductors.
  \item \textbf{First Order Circuits} \hfill \textit{(11 hrs)}\\
        First order linear differential equations with constant coefficients; simple LTI RC
        circuit; simple LTI RL circuit; analysis of miscellaneous LTI first order circuits;
        piecewise-linear first order circuits.
\end{enumerate}

\vspace{-0.2em}
{\footnotesize A more detailed course outline is available on \href{\odtuclass}{ODTÜCLASS}.}

%==============================================================
\section{Laboratory Outline}

\begin{center}
\small
\begin{tabular}{l L{5.4cm} @{\hskip 1.6em} l L{5.4cm}}
\toprule
\textbf{Exp.} & \textbf{Title} & \textbf{Exp.} & \textbf{Title}\\
\midrule
\#1 & Introduction to Voltage, Current and Resistance Measurements & \#4 & Resistive Circuits\\
\#2 & Signal Generators and Oscilloscopes   & \#5 & Operational Amplifiers\\
\#3 & LTspice Workshop \textit{(computer)}  & \#6 & First Order Circuits\\
\bottomrule
\end{tabular}
\end{center}

%==============================================================
\vspace{0.2em}
{\color{metugray}\hrule height 0.4pt}
\vspace{0.3em}
{\scriptsize\color{metugray}%
\coursecode\ \coursename\ --- \term\ --- \mysection\ (\myname).
Common syllabus prepared by the \coursecode\ teaching team;
section-specific material at \href{\notesrepo}{\notesrepotext}.\par}

\end{document}
